% !TeX root = ../../main.tex
\subsection{\bilingual{Existence for Leftmost Asymptotic Speed in \texorpdfstring{$N$}{N}-BRW}{\texorpdfstring{$N$}{N}-分支随机游走最左侧渐近速率的存在性}}

\bilingual{
    This subsection extends the result of B\'erard and Gour\'er\'e~\mycite{bérard2010brunetderridabehaviorbranchingselectionparticle} on the existence of the leftmost asymptotic speed for binary $N$-BRW to the case of infinite-branching $N$-BRW.
}{
    本小节把B\'erard与Gour\'er\'e~\mycite{bérard2010brunetderridabehaviorbranchingselectionparticle}关于二分支$N$-分支随机游走的渐近速率存在性的结论推广到无限分支$N$-分支随机游走的情形.
}

\begin{lemma}
    \label{proposition:2}
    \bilingual{
        Under \Cref{assumption:psi-one,assumption:moment-1} (i.e. $\psi(1)=0$ and $\mathbb E[(\Xi_1)_+]<\infty$), there exists $v_N=v_N(\Xi)\in\mathbb R$ such that, with probability one, and in $\mathrm L^1(\mathbb P)$, for the $N$-BRW $\mathcal V^N$ with initial population $\mathcal Z^N_0=N\delta_0$,
    }{
        在\Cref{assumption:psi-one,assumption:moment-1} (即 $\psi(1)=0$ 且 $\mathbb E[(\Xi_1)_+]<\infty$) 的条件下, 存在 $v_N=v_N(\Xi)\in\mathbb R$ 使得对于以 $\mathcal Z^N_0=N\delta_0$ 为初始粒子分布的 $N$-分支随机游走 $\mathcal V^N$, 几乎必然成立, 且在 $\mathrm L^1(\mathbb P)$ 意义下成立,
    }
    {
        \setlength{\belowdisplayskip}{-10pt}
        \begin{equation*}
            \lim_{n\to\infty}
            \frac 1n
            \mathcal M^N_n
            =
            v_N.
        \end{equation*}
    }
\end{lemma}

\begin{proof}
    \quad

    \bilingual{
        We are to apply Kingman's subadditive ergodic theorem. Consider an i.i.d.~family on the same probability space of random variables $(\Xi_{\ell,i,j})$ for $\ell\ge 0$, $i\in[\![1,N]\!]$ and $j\ge 1$ so that $\{(\Xi_{\ell,i,j})_{j\ge 1}\}_{\ell,i}$ have common distribution $\Xi$. The index $\ell$ will be used to shift the first branch from $\Xi_{0,\cdot}$ to $\Xi_{\ell,\cdot}$. Define $(W_{\ell,k})_{k\ge 0}$ for $\ell\ge 0$ as follows.
    }{
        Kingman次可加遍历定理将被应用以证明本引理. 考虑在同一概率空间上具有共同分布 $\Xi$ 的独立同分布随机变量族 $(\Xi_{\ell,i,j})_{\ell,i}$ 其中 $\ell\ge 0$, $i\in[\![1,N]\!]$ 且 $j\ge 1$. 索引 $\ell$ 将被用来把第$\ell$个对象的第一次分支的分支律从 $\Xi_{0,\cdot}$ 沿时间轴正方向移动到 $\Xi_{\ell,\cdot}$. 对 $\ell\ge 0$ 定义 $(W_{\ell,k})_{k\ge 0}$ 如下.
    }
    \begin{itemize}
        \item
        \bilingual{
            For $k=0$, let $W_{\ell,0}=N\delta_0$ for all $\ell\ge 0$.
        }{
            对 $k=0$, 对所有 $\ell\ge 0$ 记 $W_{\ell,0}=N\delta_0$.
        }
        \item
        \bilingual{
            For $k\ge 1$, suppose $W_{\ell,k-1}$ has been defined for all $\ell\ge 0$. We write
        }{
            对 $k\ge 1$, 假设 $W_{\ell,k-1}$ 已对所有 $\ell\ge 0$ 定义. 记
        }
        \begin{equation*}
            W_{\ell,k-1}=\sum_{i=1}^N\delta_{x_i},
        \end{equation*}
        \bilingual{
            with $x_1\le x_2\le\cdots\le x_N$. The population $T_{\ell,k}$ derived from $W_{\ell,k-1}$ at time $k$ by branching is
        }{
            其中 $x_1\le x_2\le\cdots\le x_N$. 由 $W_{\ell,k-1}$ 在时刻 $k$ 通过分支得到的粒子总数 $T_{\ell,k}$ 为
        }
        \begin{equation*}
            T_{\ell,k}
            :=
            \sum_{i=1}^N
            \sum_{j=1}^\infty
            \delta_{x_i+\Xi_{\ell+k,i,j}}.
        \end{equation*}
        \bilingual{
            Then $W_{\ell,k}$ is defined as being performed by the $N$ leftmost particles in $T_{\ell,k}$.
        }{
            然后 $W_{\ell,k}$ 定义为从 $T_{\ell,k}$ 中选取最左侧的 $N$ 个粒子组成的测度.
        }
    \end{itemize}
    \bilingual{
        We apply an improved version (Liggett \mycite{liggett1985improved}) of Kingman's subadditive ergodic theorem to the sequence of random variables
    }{
        Kingman次可加遍历定理的一个改进版本 (Liggett \mycite{liggett1985improved}) 将被应用于以下随机变量序列:
    }
    \begin{equation*}
        X_{m,n}:=-\min W_{m,n-m},\quad 0\le m<n.
    \end{equation*}
    \bilingual{
        We verify the assumptions of the theorem as follows.
    }{
        以下是对定理假设的验证.
    }
    \begin{itemize}
        \item (\bilingual{Subadditivity}{次可加性}).
        \bilingual{
            $X_{0,n}\le X_{0,m}+X_{m,n}$ for all $0\le m<n$ by the definition of the branching-selection procedure. Indeed,
        }{
            对所有 $0\le m<n$, 都有 $X_{0,n}\le X_{0,m}+X_{m,n}$. 事实上, 根据分支-选择步骤的定义,
        }
        \begin{equation*}
            \min W_{0,n}
            \ge
            \min W_{0,m}+\min W_{m,n-m}
        \end{equation*}
        \bilingual{
            holds since it compares the minimum of two populations of particles at time $n$ obtained by performing the branching-selection procedure specified by \Cref{corollary:2} starting from the initial population distribution at time $m$ of $W_{0,m}$ and $N\delta_{\min W_{0,m}}$, respectively.
        }{
            成立, 此因不等号两侧比较了两种按照\Cref{corollary:2} 具体说明的分支-选择步骤所得到的时刻 $n$ 的粒子总数的最小值, 它们在时刻 $m$ 初始粒子分布分别为 $W_{0,m}$ 和 $N\delta_{\min W_{0,m}}$.
        }
        \item (\bilingual{Stationarity}{平稳性}).
        \bilingual{
            $(X_{nk,(n+1)k})_{n\ge 1}$ is a stationary sequence for each $k\ge 1$. Indeed, $X_{nk,(n+1)k}=-\min W_{nk,k}$. By definition of $W$, for each $k\ge 1$, the sequence $(W_{nk,k})_{n\ge 1}$ is i.i.d., since, for any $1\le n_1<n_2$, we have the fact that $n_1k+k-1<n_2k$ and that $W_{n_1k,k}$ is constructed from the point processes
        }{
            对任一 $k\ge 1$, $(X_{nk,(n+1)k})_{n\ge 1}$ 是一个平稳序列. 事实上, 回顾 $X$ 的定义有 $X_{nk,(n+1)k}=-\min W_{nk,k}$. 根据 $W$ 的定义, 对每个 $k\ge 1$, 序列 $(W_{nk,k})_{n\ge 1}$ 是独立同分布的, 此因对于任意 $1\le n_1<n_2$, 都有 $n_1k+k-1<n_2k$ 且 $W_{n_1k,k}$ 是由如下点过程构建的
        }
        $$
            \left\{
                \left(\Xi_{n_1k+l,i,j}\right)_{j\ge 1}:
                l\in[\![0,k-1]\!],\,
                i\in[\![0,N]\!]
            \right\},
        $$
        \bilingual{
            whereas $W_{n_2k,k}$ is constructed in the same way from the point processes
        }{
            然而 $W_{n_2k,k}$ 是以同样的方式由如下点过程构建的
        }
        $$
            \left\{
                \left(\Xi_{n_2k+l,i,j}\right)_{j\ge 1}:
                l\in[\![0,k-1]\!],\,
                i\in[\![0,N]\!]
            \right\},
        $$
        \bilingual{
            and these two sets of point processes are independently identically distributed. Hence, the sequence $(X_{nk,(n+1)k})_{n\ge 1}$ is also i.i.d. The result then follows from the fact that i.i.d.~sequences are stationary.
        }{
            且这两组点过程是独立同分布的. 从而, 序列 $(X_{nk,(n+1)k})_{n\ge 1}$ 也是独立同分布的. 结果由独立同分布序列是平稳的这一事实得出.
        }
        \item (\bilingual{Ergodicity}{遍历性}).
        \bilingual{
            Moreover, $(X_{nk,(n+1)k})_{n\ge 1}$ is an ergodic sequence for each $k\ge 1$. Indeed, since the sequence is i.i.d. as stated in the last item, the result follows from the fact that i.i.d.~sequences are ergodic.
        }{
            进一步, 对任一 $k\ge 1$, $(X_{nk,(n+1)k})_{n\ge 1}$ 是一个遍历序列. 事实上, 因为该序列是独立同分布的 (如上一项所述), 结果由独立同分布序列是遍历的这一事实得出.
        }
        \item (\bilingual{Distribution invariance}{分布不变性}).
        \bilingual{
            The distribution of $\{X_{m,m+k}:k\ge 1\}$ does not depend on $m$. Indeed, $X_{m,m+k}=-\min W_{m,k}$, and by definition of $W$, for any $0\le m_1<m_2$, $W_{m_1,k}$ and $W_{m_2,k}$ are identically distributed, since $W_{m_1,k}$ is constructed from the point processes
        }{
            $\{X_{m,m+k}:k\ge 1\}$ 的分布不依赖于 $m$. 事实上, 回顾 $X$ 的定义有 $X_{m,m+k}=-\min W_{m,k}$, 根据 $W$ 的定义, 对任意 $0\le m_1<m_2$, $W_{m_1,k}$ 和 $W_{m_2,k}$ 是同分布的, 此因 $W_{m_1,k}$ 是由如下点过程构建的
        }
        $$
            \left\{
                \left(\Xi_{m_1+l,i,j}\right)_{j\ge 1}:
                l\in[\![0,k-1]\!],\,
                i\in[\![0,N]\!]
            \right\},
        $$
        \bilingual{
            while $W_{m_2,k}$ is constructed in the same way from the point processes
        }{
            同时 $W_{m_2,k}$ 是以同样的方式由如下点过程构建的
        }
        $$
            \left\{
                \left(\Xi_{m_2+l,i,j}\right)_{j\ge 1}:
                l\in[\![0,k-1]\!],\,
                i\in[\![0,N]\!]
            \right\},
        $$
        \bilingual{
            and these two sets of point processes are identically distributed. Hence, the distribution of $\{X_{m,m+k}:k\ge 1\}$ does not depend on $m$.
        }{
            且这两组点过程是同分布的. 从而, $\{X_{m,m+k}:k\ge 1\}$ 的分布不依赖于 $m$.
        }
        \item (\bilingual{Finite expectation, Part 1}{有限期望, 部分 1}).
        \bilingual{
            For each $n\ge 1$, $\mathbb E[X_{0,n}]\ge cn$ for some
            $c> -\infty$. Follow recursively one currently selected particle
            and, before the next selection, its leftmost child. If $R_k$
            denotes the displacement used at step $k$, then
            $(R_k)_{k\ge0}$ are independent copies of $\Xi_1$ and selection
            can only move the population minimum further to the left.
            Consequently,
        }{
            对每个 $n\ge 1$, $\mathbb E[X_{0,n}]\ge cn$ 对某个
            $c>-\infty$ 成立. 递归地选取当前一个被保留的粒子, 并在下一次选择前追踪其最左子代. 若第 $k$ 步采用的位移记为 $R_k$, 则
            $(R_k)_{k\ge0}$ 是 $\Xi_1$ 的独立同分布副本, 且选择操作只可能使种群最小位置进一步向左. 因此
        }
        \begin{equation*}
            \min W_{0,n}\le\sum_{k=0}^{n-1}R_k,
            \qquad
            X_{0,n}\ge-\sum_{k=0}^{n-1}R_k.
        \end{equation*}
        \bilingual{
            Under \Cref{assumption:psi-one},
            $\mathbb E[(\Xi_1)_-]<\infty$ by
            \Cref{lemma:moment-1:neg}; together with
            \Cref{assumption:moment-1}, $\Xi_1$ is integrable. Hence
            $\mathbb E[X_{0,n}]\ge-n\mathbb E[\Xi_1]$.
        }{
            在\Cref{assumption:psi-one} 下, 由
            \Cref{lemma:moment-1:neg} 有
            $\mathbb E[(\Xi_1)_-]<\infty$; 再结合
            \Cref{assumption:moment-1}, $\Xi_1$ 可积. 因而
            $\mathbb E[X_{0,n}]\ge-n\mathbb E[\Xi_1]$.
        }
        \item (\bilingual{Finite expectation, Part 2}{有限期望, 部分 2}).
        \bilingual{
            For each $n\ge1$, $\mathbb E[|X_{0,n}|]<\infty$.
            Subadditivity gives
        }{
            对每个 $n\ge1$, $\mathbb E[|X_{0,n}|]<\infty$. 次可加性给出
        }
        \begin{equation*}
            (X_{0,n})_+
            \le
            \sum_{k=1}^n(X_{k-1,k})_+.
        \end{equation*}
        \bilingual{
            The designated-line inequality from Part 1 gives the other side,
        }{
            部分 1 中的指定谱系不等式给出另一侧,
        }
        \begin{equation*}
            (X_{0,n})_-
            \le
            \left(\sum_{k=0}^{n-1}R_k\right)_+
            \le
            \sum_{k=0}^{n-1}(R_k)_+.
        \end{equation*}
        \bilingual{
            For the one-step positive part,
        }{
            对于一步正部,
        }
        \begin{align*}
            (X_{0,1})_+
            &=\left(\min_{i\in[\![1,N]\!]}\Xi_{0,i,1}\right)_-
            =\max_{i\in[\![1,N]\!]}(\Xi_{0,i,1})_-
            \le\sum_{i=1}^N(\Xi_{0,i,1})_-.
        \end{align*}
        \bilingual{
            Thus $\mathbb E[(X_{0,1})_+]\le
            N\mathbb E[(\Xi_1)_-]<\infty$ by
            \Cref{lemma:moment-1:neg}, while
            $\mathbb E[(X_{0,n})_-]\le
            n\mathbb E[(\Xi_1)_+]<\infty$ by
            \Cref{assumption:moment-1}. Stationarity and the positive-part
            bound now prove $\mathbb E[|X_{0,n}|]<\infty$. This argument does
            not use an expectation identity before integrability is known.
        }{
            因此由\Cref{lemma:moment-1:neg} 有
            $\mathbb E[(X_{0,1})_+]\le
            N\mathbb E[(\Xi_1)_-]<\infty$; 同时由
            \Cref{assumption:moment-1} 有
            $\mathbb E[(X_{0,n})_-]\le
            n\mathbb E[(\Xi_1)_+]<\infty$. 再结合平稳性与正部上界, 即得
            $\mathbb E[|X_{0,n}|]<\infty$. 此论证没有在可积性得到证明之前使用期望恒等式.
        }
    \end{itemize}
    \bilingual{
        We conclude that the hypotheses of Kingman's subadditive ergodic theorem are satisfied, and thus
    }{
        定理的假设已被验证, 因此 Kingman 次可加遍历定理适用, 从而
    }
    \begin{equation*}
        -v_N
        :=
        \lim_{n\to\infty}
        \frac 1n
        X_{0,n}
        =
        -\lim_{n\to\infty}
        \frac 1n
        \min W_{0,n}
    \end{equation*}
    \bilingual{
        exists both almost surely and in $\mathrm L^1(\mathbb P)$, and is constant. By the definition of $W$, $\mathcal M^N_n$ has the same distribution as $W_{0,n}$ for all $n\ge 0$, and thus the desired result follows.
    }{
        在几乎必然的意义下与在 $\mathrm L^1(\mathbb P)$ 意义下都存在, 且是常数. 根据 $W$ 的定义, 所有 $n\ge 0$, $\mathcal M^N_n$ 与 $W_{0,n}$ 同分布, 从而命题得证.
    }
\end{proof}

\bilingual{
    We extend the last result to more general initial population distributions by using the coupling results in the last subsection.
}{
    利用上一小节的耦合结论, 上述结果可被推广到更一般的初始粒子分布.
}
\begin{corollary}
    \label{corollary:leftmost-speed:N}
    \bilingual{
        Under \Cref{assumption:psi-one,assumption:moment-1} (i.e. $\psi(1)=0$ and $\mathbb E[(\Xi_1)_+]<\infty$), there exists $v_N=v_N(\Xi)\in\mathbb R$ such that, with probability one, and in $\mathrm L^1(\mathbb P)$, for the $N$-BRW $\mathcal V^N$ with initial population $\mathcal Z^N_0\in \mathcal C(N)$,
    }{
        在\Cref{assumption:psi-one,assumption:moment-1} (即 $\psi(1)=0$ 且 $\mathbb E[(\Xi_1)_+]<\infty$) 的条件下, 存在 $v_N=v_N(\Xi)\in\mathbb R$ 使得对于以任一 $\mathcal Z^N_0\in \mathcal C(N)$ 为初始粒子分布的 $N$-分支随机游走 $\mathcal V^N$, 几乎必然成立, 且在 $\mathrm L^1(\mathbb P)$ 意义下成立,
    }
    {
        \setlength{\belowdisplayskip}{-10pt}
        \begin{equation*}
            \lim_{n\to\infty}
            \frac 1n
            \mathcal M^N_n
            =
            v_N.
        \end{equation*}
    }
\end{corollary}

\begin{proof}
    \quad

    \bilingual{
        Note that
    }{
        注意到
    }
    \begin{equation*}
        \mathcal Z'_0
        :=
        N\delta_{\max\mathcal Z^N_0}
        \succcurlyeq
        \mathcal Z^N_0
        \succcurlyeq
        N\delta_{\min\mathcal Z^N_0}
        =:
        \mathcal Z''_0.
    \end{equation*}
    \bilingual{
        By \Cref{corollary:2} and \Cref{corollary:translation}, for another $N$-BRW $\mathcal V$ with initial population distribution $N\delta_0$, there exists a coupling of all these four $N$-BRWs so that
    }{
        根据\Cref{corollary:2} 与\Cref{corollary:translation}, 对另一个以 $N\delta_0$ 为初始粒子分布的 $N$-分支随机游走 $\mathcal V$, 存在一个耦合使得这四个 $N$-分支随机游走满足如下关系:
    }
    \begin{equation*}
        \max\mathcal Z_0^N
        +
        \min\mathcal V_n
        =
        \min \mathcal Z'_n
        \ge \min\mathcal Z^N_n
        \ge \min\mathcal Z''_n
        =
        \min\mathcal Z_0^N
        +
        \min\mathcal V_n,
    \end{equation*}
    \bilingual{
        almost surely for all $n\ge 0$. By \Cref{proposition:2},
        there exists $v_N=v_N(\Xi)\in\mathbb R$ such that, with probability one, and in $\mathrm L^1(\mathbb P)$, for the $N$-BRW $\mathcal V$ with initial population $N\delta_0$, $\lim_{n\to\infty}\frac 1n\min\mathcal V_n=v_N$. It immediately follows that
    }{
        对所有 $n\ge 0$ 几乎必然成立. 根据\Cref{proposition:2}, 存在 $v_N=v_N(\Xi)\in\mathbb R$ 使得对于以 $N\delta_0$ 为初始粒子分布的 $N$-分支随机游走 $\mathcal V$, 几乎必然成立, 且在 $\mathrm L^1(\mathbb P)$ 意义下成立 $\lim_{n\to\infty}\frac 1n\min\mathcal V_n=v_N$. 由此立即得到
    }
    \begin{equation*}
        \lim_{n\to\infty}
        \frac 1n\min\mathcal Z'_n
        =
        \lim_{n\to\infty}
        \frac 1n\min\mathcal Z''_n
        =
        \lim_{n\to\infty}
        \frac 1n\min\mathcal Z^N_n
        =
        v_N
    \end{equation*}
    \bilingual{
        almost surely and in $\mathrm L^1(\mathbb P)$, which completes the proof.
    }{
        几乎必然成立, 且在 $\mathrm L^1(\mathbb P)$ 意义下成立, 由此命题得证.
    }
\end{proof}

\bilingual{
    We further extend the last result to initial population distributions with total mass less than $N$.
}{
    我们进一步将上述结果推广到总质量小于 $N$ 的初始粒子分布.
}

\begin{corollary}
    \label{corollary:leftmost-speed:N0}
    \bilingual{
        Consider \Cref{assumption:psi-one,assumption:moment-1} (i.e. $\psi(1)=0$ and $\mathbb E[(\Xi_1)_+]<\infty$) and an additional \Cref{assumption:psi-zero} (i.e., $\psi(0)>0$). There exists $v_N=v_N(\Xi)\in\mathbb R$ such that, with probability one, and in $\mathrm L^1(\mathbb P)$, for the $N$-BRW $\mathcal V^N$ with initial population $\mathcal Z^N_0\in \mathcal C(N_0)$ and $N_0\le N$,
    }{
        考虑\Cref{assumption:psi-one,assumption:moment-1} (即 $\psi(1)=0$ 且 $\mathbb E[(\Xi_1)_+]<\infty$) 和额外的\Cref{assumption:psi-zero} (即 $\psi(0)>0$). 存在 $v_N=v_N(\Xi)\in\mathbb R$ 使得对于具初始粒子分布 $\mathcal Z^N_0\in \mathcal C(N_0)$ 且 $N_0\le N$ 的 $N$-分支随机游走 $\mathcal V^N$,
    }
    {
        \setlength{\belowdisplayskip}{-10pt}
        \begin{equation*}
            \lim_{n\to\infty}
            \frac 1n
            \mathcal M^N_n
            =
            v_N.
        \end{equation*}
    }
\end{corollary}

\begin{proof}
    \bilingual{
        Under \Cref{assumption:psi-zero}, i.e., $\psi(0)>0$ or $\mathbb E[\#\Xi]>0$, we have $\mathbb P(\#\Xi\ge 2)>0$. Then there exists some time $n_0$ such that $\#\mathcal Z^N_{n_0}=N$. Then the result follows from \Cref{corollary:leftmost-speed:N}.
    }{
        在\Cref{assumption:psi-zero} 下, 即 $\psi(0)>0$ 或 $\mathbb E[\#\Xi]>0$, 我们有 $\mathbb P(\#\Xi\ge 2)>0$. 于是存在某个时间 $n_0$ 使得 $\#\mathcal Z^N_{n_0}=N$. 由此结果可直接从\Cref{corollary:leftmost-speed:N} 得出.
    }
\end{proof}
